\chapter{Introduction}
\label{chap:introduction}

\section{Background}

Drug discovery is a long and expensive process in which a large number of
candidate molecules must be screened before a small subset can proceed to
experimental and clinical validation. One of the central questions in this
process is whether a drug molecule can interact with a biological target, such
as a protein, enzyme, receptor, or gene product. These drug--target
interactions (DTIs) provide important evidence for understanding therapeutic
mechanisms, identifying potential side effects, and discovering new uses for
existing drugs.

Traditional experimental methods for identifying DTIs are reliable but costly.
They require domain expertise, laboratory resources, and substantial time. As a
result, computational DTI prediction has become an important research direction
in bioinformatics and machine learning. Instead of testing every possible
drug--target pair experimentally, computational models estimate which pairs are
most likely to interact. These predictions can then be used to prioritize
candidate interactions for further biological validation.

Early computational approaches often represented the DTI problem as a
relationship between two spaces: the chemical space of drugs and the genomic or
biological space of target proteins. For example, the benchmark work of
Yamanishi et al. modeled DTI prediction by integrating chemical similarity and
genomic similarity information \cite{yamanishi2008prediction}. This view is
still influential because it frames DTI prediction as a supervised learning
problem over pairs of heterogeneous entities.

However, modern biomedical data is no longer limited to pairwise drug and
protein descriptors. Drugs, proteins, diseases, pathways, side effects, and
biological processes are connected through many relation types. These
relationships naturally form a biomedical knowledge graph. A knowledge graph
can provide context that is difficult to capture from isolated descriptors
alone. For example, two drugs may not appear similar by molecular fingerprint,
but they may be connected to similar diseases, pathways, or protein families in
the graph. Likewise, two proteins may share functional context even when their
sequence-level descriptors are not sufficient.

This project studies a machine learning pipeline that uses both descriptor
features and knowledge graph representation learning. The first part of the
pipeline learns entity representations from biomedical graph triples using
knowledge graph embedding (KGE) models. The second part combines these learned
representations with drug and protein descriptors and trains a neural
factorization machine (NFM) for final DTI classification.

\section{Problem Context}

The DTI prediction task can be formulated as a binary prediction problem. Given
a drug $d$ and a target protein $t$, the model predicts whether the pair
$(d,t)$ represents a true interaction. Let $\mathcal{D}$ denote the set of
drugs and $\mathcal{T}$ denote the set of target proteins. Each observed pair is
associated with a label:

\begin{equation}
    y_{dt} =
    \begin{cases}
        1, & \text{if drug } d \text{ is known to interact with target } t, \\
        0, & \text{otherwise.}
    \end{cases}
\end{equation}

The objective is to learn a scoring function:

\begin{equation}
    f_{\theta}: \mathcal{D} \times \mathcal{T} \rightarrow [0,1],
\end{equation}

where $f_{\theta}(d,t)$ estimates the probability that drug $d$ interacts with
target $t$. A higher score indicates that the model considers the pair more
likely to be a valid interaction.

Although this formulation is simple, the task is challenging in practice. First,
known DTI data is sparse compared with the number of possible drug--target
pairs. Second, many unobserved pairs are not guaranteed to be true negatives;
they may simply be unknown interactions. Third, the model must combine
heterogeneous information sources, including chemical features, protein
features, and relational graph structure. Finally, the task becomes harder under
cold-start settings, where the model must predict interactions for drugs or
proteins that have limited observed interaction history.

\section{Motivation}

The main motivation of this project is to improve DTI prediction by combining
two complementary sources of information: descriptor-based features and
knowledge graph embeddings.

Descriptor features provide direct information about the entities. Drug
descriptors may encode molecular structure, chemical fingerprints, or other
compound-level properties. Protein descriptors may encode sequence-derived
features or biological properties. These features are useful because they give
the supervised classifier explicit information about the drug and target in each
candidate pair.

Knowledge graph embeddings provide relational information. KGE models learn
dense vector representations of entities and relations from graph triples of
the form $(h,r,t)$, where $h$ is a head entity, $r$ is a relation type, and $t$
is a tail entity. A simple model such as DistMult can efficiently learn entity
and relation embeddings through a bilinear scoring function
\cite{yang2015embedding}. However, DistMult has limited expressiveness because
its scoring function is symmetric and does not explicitly aggregate
neighborhood structure.

More expressive graph representation models, such as CompGCN, address part of
this limitation by performing relation-aware message passing over
multi-relational graphs \cite{vashishth2020composition}. This is especially
relevant for biomedical data because a drug or protein is often better
understood through its surrounding graph context rather than through isolated
pairwise relations.

After graph representations are obtained, the remaining challenge is how to use
them effectively in the supervised prediction stage. Neural factorization
machines are designed to model feature interactions in sparse predictive
settings \cite{he2017neural}. This makes NFM a suitable candidate for combining
drug identifiers, target identifiers, descriptor features, and graph-derived
embeddings. In this project, the NFM component is used as the final classifier
that learns interaction patterns from the constructed pair-level features.

\section{Research Gap}

Many DTI prediction methods rely mainly on handcrafted descriptors or similarity
matrices. These methods can be effective when high-quality descriptors are
available, but they may fail to exploit broader relational structure. On the
other hand, graph-based methods can capture biomedical context, but graph
embeddings alone may not fully represent the chemical and biological properties
of individual drugs and proteins.

This project is motivated by the gap between these two perspectives. A
descriptor-only model may miss relational evidence, while a graph-only model may
ignore useful entity-level features. Therefore, a hybrid pipeline is explored:
KGE models learn representations from biomedical knowledge graphs, and an NFM
model combines those representations with descriptor features for supervised DTI
prediction.

The project also compares different KGE choices. DistMult is treated as a
lightweight baseline because it is efficient and widely used. CompGCN is
considered as a stronger relational graph encoder because it can incorporate
neighborhood information and jointly update entity and relation representations.
This comparison helps clarify whether more expressive graph representation
learning improves the downstream DTI prediction task.

\section{Objectives}

The main objective of this project is to develop and evaluate a reproducible
machine learning pipeline for drug--target interaction prediction using
knowledge graph embeddings and neural factorization machines.

The specific objectives are:

\begin{itemize}
    \item to study DTI prediction as a supervised binary classification problem;

    \item to preprocess benchmark DTI datasets and biomedical knowledge graph
    triples for model training;

    \item to train KGE models, including DistMult and CompGCN, for learning
    drug and protein representations;

    \item to construct pair-level features by combining descriptor features and
    graph-derived embeddings;

    \item to train an NFM classifier for final interaction prediction;

    \item to evaluate the model using ROC-AUC and PR-AUC across selected
    datasets and data splits;

    \item to compare the proposed KGE-NFM pipeline with relevant baseline
    approaches.
\end{itemize}

\section{Scope of the Project}

This project focuses on computational prediction rather than experimental
validation. The model outputs probability scores for candidate drug--target
pairs, but these predictions should be interpreted as prioritization signals
rather than confirmed biological interactions.

The current implementation supports datasets such as \textit{Yamanishi08},
\textit{BioKG}, and \textit{Hetionet}. The main pipeline is implemented in
Python using PyTorch for neural model training and PyKEEN for knowledge graph
embedding models.

\begin{description}
    \item[In scope:] dataset loading, preprocessing, KGE training, entity
    embedding extraction, feature construction, NFM training, baseline
    comparison, and metric reporting.

    \item[Out of scope:] wet-lab validation, clinical validation, deployment of
    a production drug discovery platform, and detailed biological mechanism
    interpretation of each predicted interaction.
\end{description}

\section{Expected Contributions}

The expected contributions of this project are:

\begin{itemize}
    \item a clear implementation of a KGE-NFM pipeline for DTI prediction;

    \item an experimental comparison between lightweight and more expressive KGE
    models in the downstream DTI setting;

    \item an evaluation framework using fold-based prediction outputs, ROC
    curves, PR curves, and aggregate AUC metrics;

    \item a reusable report and code structure for future experiments on
    biomedical graph-based prediction tasks.
\end{itemize}

\section{Report Organization}

The rest of this report is organized as follows.

\begin{table}[H]
\centering
\caption{Report organization}
\label{tab:report_organization}
\begin{tabular}{cl}
\toprule
\textbf{Chapter} & \textbf{Main content} \\
\midrule
Chapter 1 & Background, motivation, objectives, scope, and contributions \\
Chapter 2 & DTI, biomedical knowledge graph, and KGE background \\
Chapter 3 & Yamanishi08 data analysis and warm-start settings \\
Chapter 4 & KGE-NFM methodology with DistMult and CompGCN encoders \\
Chapter 5 & Experimental setup, results, and discussion \\
Chapter 6 & Conclusion, limitations, and future work \\
\bottomrule
\end{tabular}
\end{table}

